%% Owner: Bernat Ventura
\section{Stack: app}
\label{sec:stack-app}

The app stack is a ratings monitor. The statement suggests, as the
automatic feeder, ``a monitor that captures comment ratings on web
pages'' \parencite{llusa2026project}, and the whole application is
built around that idea, so every service has a real job. Many web pages
publish their rating as schema.org JSON-LD (\texttt{aggregateRating}),
the format search engines read to show stars. The feeder reads that
rating from a list of pages, the core application stores it and shows
the latest rating of each page at
\url{https://app.ccds-project.duckdns.org}, and it can email the
summary as a report. The pages are two books on Goodreads, two recipes
on BBC Good Food and a film on Rotten Tomatoes. Other sites, such as
IMDb or Allrecipes, were tested and dropped because they do not publish
the rating in that format.

\subsection{Services and images}

\Cref{tab:app-services} lists the seven services. Five images are built
by the team and pushed to the registry of the basic stack as
\texttt{ccds-project.duckdns.org/app-<service>:<version>}; Redis uses the
official image unchanged.

\begin{table}[htbp]
  \centering
  \caption{Services of the app stack.}
  \label{tab:app-services}
  \small
  \begin{tabular}{llll}
    \toprule
    Service & Image                         & Replicas & Role \\
    \midrule
    core    & \texttt{app-core} (Flask)     & 2 & web page and API \\
    static  & \texttt{app-static} (httpd)   & 2 & static files under \texttt{/static/} \\
    db      & \texttt{app-db} (PostgreSQL 17) & 1 & ratings, data on NFS \\
    redis   & \texttt{redis:8.2-alpine}     & 1 & 30\,s cache of the summary \\
    smtp    & \texttt{app-smtp} (msmtpd)    & 2 & relays mail to Gmail on port 587 \\
    backup  & \texttt{app-backup} (crond)   & 1 & nightly \texttt{pg\_dump} to NFS \\
    auto    & \texttt{app-auto} (Python)    & 2 & reads the ratings, posts them to core \\
    \bottomrule
  \end{tabular}
\end{table}

A rating flows through the stack as follows. Each \texttt{auto} replica
downloads its pages, extracts the \texttt{aggregateRating} with a
regular expression and a search of the JSON tree, and posts it to
\texttt{core} with a shared token. \texttt{core} inserts it into one
PostgreSQL table. The summary is a single
\texttt{SELECT DISTINCT ON (url)} that keeps the latest rating of each
page. It follows the cache-aside pattern: \texttt{core} first asks
Redis, and on a miss queries PostgreSQL and stores the result for
30\,seconds. If Redis is down, the page still works from the database.
\texttt{POST /api/report} sends the summary by email through
\texttt{smtp}.

\subsection{Dockerfiles}

\begin{description}
  \item[core and auto] \texttt{python:3.13-slim}, pinned dependencies,
    and an unprivileged user. \texttt{core} runs gunicorn with one
    worker and eight threads: the stack scales by adding replicas, not
    workers inside a container.
  \item[static] \texttt{httpd:2.4-alpine} with the files copied under
    \texttt{htdocs/static/}, because nginx forwards the path unchanged,
    plus cache headers.
  \item[db] The official PostgreSQL image with a small entrypoint in
    front. It creates the data directory on NFS, owned by
    \texttt{postgres} (uid 70), and then runs the official entrypoint
    as \texttt{postgres}. \Cref{sec:app-nfs} explains why.
  \item[smtp] \texttt{alpine:3.22} with \texttt{msmtpd}, a minimal SMTP
    server that hands every message to \texttt{msmtp}, which relays it
    to the provider with authentication and STARTTLS. Azure blocks
    outbound port 25 \parencite{azure-smtp}, so the relay uses port 587.
  \item[backup] The same PostgreSQL image, so \texttt{pg\_dump} matches
    the server version, with busybox \texttt{crond}.
\end{description}

\subsection{Local and production files}

\texttt{compose.yml} runs the stack on a laptop with
\texttt{docker compose up}. It is self-contained: it adds its own nginx
on port 8080 in place of the basic stack, replaces the relay with
Mailpit, which catches every email and shows it on port 8025, keeps the
data in named volumes, passes passwords as plain environment variables,
and feeds two demo pages served by \texttt{static}, so it works without
internet access. \texttt{production.yml} is the Swarm version: images
from the registry, data on NFS, Swarm secrets, the external
\texttt{proxy} network shared with nginx, replicas and update policies.

\subsection{Replicas, placement and updates}

The four stateless services run two replicas each, and Swarm may place
them on any node. Three services run one replica, all within the
exceptions of the statement \parencite{llusa2026project}:
\texttt{db} is a DBMS, \texttt{backup} is a cron job, and \texttt{redis}
is an in-memory database used as a shared cache, since two independent
instances behind one virtual IP would hold two different caches.

The feeder replicas share the work without talking to each other. Swarm
gives every task a slot number, \texttt{\{\{.Task.Slot\}\}}, passed as
\texttt{FEEDER\_SLOT}. With $N$ replicas, replica $k$ takes URLs
$k, k+N, k+2N, \ldots$ (\texttt{URLS[SLOT-1::REPLICAS]} in Python), so
each page is read by exactly one replica.

Stateless services update one task at a time, start the new task before
stopping the old one, and roll back if the update fails. The database
uses \texttt{stop-first} instead, because two PostgreSQL servers must
never open the same data directory.

\subsection{Data, NFS and secrets}
\label{sec:app-nfs}

PostgreSQL keeps its data in \texttt{/srv/nfs/app/db} and the dumps go
to \texttt{/srv/nfs/app/backups}, both on the NFS share that every node
mounts. Swarm refuses a bind mount whose source directory does not
exist, so \texttt{db} and \texttt{backup} mount the whole share and
create their own directory on start. This only works if root in a
container can write to the share. The export first used
\texttt{root\_squash}, which maps root to \texttt{nobody} on the client
nodes. That broke the registry, which ran on manager-2 and could not
store new images, and it would have broken the database too. The export
now uses \texttt{no\_root\_squash}, restricted to the cluster subnet
\texttt{10.0.1.0/24}.

The data really is shared. During a redeploy Swarm stopped PostgreSQL on
manager-1 and started it on manager-3, and the new server found the
data where the old one had left it.

Three values are Swarm secrets: the database password, the token that
protects the write endpoints, and the SMTP password, which is a Gmail
app password of an account created for the project. Their files live
only on the deploying node, in \texttt{app/secrets/}, which git ignores
because the repository is public. Secrets cannot change in place, so
each name carries a version (\texttt{app\_smtp\_password\_v2}) that is
raised when the file changes.

The dumps are described with the other databases in \cref{sec:data}:
\texttt{pg\_dump} in custom format every night at 03:00, a check with
\texttt{pg\_restore --list} before a dump is kept, and seven days of
retention. Redis holds nothing worth a backup.

\subsection{Health checks and metrics}

\texttt{core} answers \texttt{/healthz} and \texttt{static} serves its
stylesheet; both are image health checks, so Swarm replaces a task that
stops answering. The \texttt{smtp} check only tests that \texttt{msmtpd}
runs: opening an SMTP session every 30\,seconds filled the logs with
aborted sessions. \texttt{core} exposes Prometheus metrics at
\texttt{/metrics}: cache hits and misses, and ratings received. nginx
answers 404 to that path from outside. Prometheus does not scrape it:
the stats stack only scrapes node exporter and cAdvisor.

\subsection{Testing}

Locally, \texttt{docker compose up --build} built the five images, the
feeder read both demo pages, \texttt{/api/pages} returned their
ratings, and the report email appeared in Mailpit. The first local run
failed with a 502: the local nginx had resolved \texttt{core} when it
started and kept the old address after \texttt{core} was recreated. It
now resolves the name on every request, as the nginx of the basic stack
does.

On the Swarm, all seven services reached their replica count across the
three nodes. The public endpoint returns the five pages, and a report
reached a real mailbox through Gmail. Two problems appeared on the
way: the registry could not store images because of
\texttt{root\_squash}, and the first SMTP attempt failed because the
user name was still the example value; the relay log showed Gmail's
answer, \texttt{535 Username and Password not accepted}.
